% Source PDF: source/普通高考/2012/2012辽宁文.pdf, page 1, question 10.
% Original q10 is vector-style flowchart; redraw keeps its node/edge topology.
% Node-edge list: 开始→S=4→i=1→i<6?;
% 是→S=2/(2-S)→i=i+1→left return to the trunk before the decision;
% 否→输出 S→结束. No other directed edges are present.
% Evidence: tmp/review_2012_liaoning_liberal/grid/page_grid100.png,
%   tmp/review_2012_liaoning_liberal/grid/page_grid50.png,
%   tmp/review_2012_liaoning_liberal/crop/q10_orig_fig_final.png.

\begin{tikzpicture}[
  baseline,
  >=Stealth,
  x=1cm,
  y=0.96cm,
  line width=0.45pt,
  line cap=round,
  line join=round,
  font=\normalsize,
  flowtext/.style={anchor=center,align=center,inner sep=0pt,
    execute at begin node={\setlength{\parindent}{0pt}}},
  flowbox/.style={flowtext,draw,minimum height=.68cm,inner xsep=3pt,inner ysep=2pt},
  branchlabel/.style={flowtext,inner sep=0pt},
  terminal/.style={flowbox,rounded corners=.18cm,minimum width=1.35cm},
  process/.style={flowbox},
  io/.style={flowbox,shape=trapezium,trapezium left angle=72,trapezium right angle=108,
    minimum height=.76cm},
  decision/.style={flowbox,shape=diamond,aspect=2.8,minimum width=3.65cm,
    minimum height=1.00cm,inner sep=0pt}
]
  \node[terminal] (start) at (0,8.70) {开始};
  \node[process,minimum width=1.70cm] (initS) at (0,7.48) {$S=4$};
  \node[process,minimum width=1.70cm] (initi) at (0,6.28) {$i=1$};
  \node[decision] (test) at (0,4.86) {$i<6$};
  \node[process,minimum width=3.00cm,minimum height=1.00cm] (update) at (0,3.38)
    {$S=\dfrac{2}{2-S}$};
  \node[process,minimum width=2.25cm] (inc) at (0,1.95) {$i=i+1$};
  \node[io,minimum width=2.05cm] (output) at (3.85,4.05) {输出 $S$};
  \node[terminal] (finish) at (3.85,2.62) {结束};

  \draw[->] (start.south) -- (initS.north);
  \draw[->] (initS.south) -- (initi.north);
  \draw[->] (initi.south) -- (test.north);
  \draw[->] (test.south) -- node[branchlabel,left=2pt] {是} (update.north);
  \draw[->] (update.south) -- (inc.north);
  \draw[->] (test.east) -- node[branchlabel,above=2pt] {否} (3.85,4.86) -- (output.north);
  \draw[->] (output.south) -- (finish.north);

  \coordinate (loop-bottom) at (0,0.72);
  \coordinate (loop-left) at (-2.45,0.72);
  \coordinate (loop-entry) at (-2.45,5.55);
  \draw (inc.south) -- (loop-bottom) -- (loop-left) -- (loop-entry);
  \draw[->] (loop-entry) -- (0,5.55);

  \path[use as bounding box] (-2.85,0.42) rectangle (4.75,9.02);
\end{tikzpicture}
